% ==============================================================================
% Appendix A: Computational Infrastructure and Visualization
% ==============================================================================

\chapter{Computational Infrastructure and Visualization}
\label{app:computational}

This appendix documents the computational infrastructure supporting the
compendium: the Python module catalog, algorithm pseudocode, time and space
complexity estimates, hardware requirements, and the interactive visualization
suite.  All paths are relative to the repository root.

% ==============================================================================
\section{Module Catalog}
\label{app:module_catalog}
% ==============================================================================

The \texttt{src/mathphysics/} package contains 29 Python modules organized by
function.  Table~\ref{tab:module_catalog} lists each module, its primary
responsibility, and key public API.

\begin{longtable}{p{4.5cm}p{4cm}p{5cm}}
\toprule
\textbf{Module} & \textbf{Responsibility} & \textbf{Key API} \\
\midrule
\endhead
\texttt{algebra.py} &
    Lie algebra computations (Cartan matrices, roots, Weyl groups) &
    \texttt{LieAlgebra}, \texttt{generate\_roots()} \\
\texttt{clifford.py} &
    Clifford algebra and spinor representations &
    \texttt{CliffordAlgebra}, \texttt{gamma\_matrices()} \\
\texttt{jordan.py} &
    Exceptional Jordan algebras ($\mathfrak{J}_3(\mathbb{O})$) &
    \texttt{JordanAlgebra}, \texttt{jordan\_triple()} \\
\texttt{kac\_moody.py} &
    Kac--Moody and affine algebra extensions &
    \texttt{KacMoodyAlgebra}, \texttt{affine\_extension()} \\
\texttt{invariant\_theory.py} &
    Polynomial invariants, quartic $\EE{7}$ invariant &
    \texttt{quartic\_e7\_invariant()}, \texttt{ring\_of\_invariants()} \\
\texttt{projective\_geometry.py} &
    Finite projective spaces PG$(n,q)$ &
    \texttt{ProjectiveSpace}, \texttt{points()}, \texttt{lines()} \\
\texttt{lattice\_theory.py} &
    Root lattices, sphere packing, theta series &
    \texttt{RootLattice}, \texttt{theta\_series()} \\
\texttt{modular\_forms.py} &
    Modular forms, $j$-invariant, eta functions &
    \texttt{j\_invariant()}, \texttt{eta\_function()} \\
\texttt{fractal\_analysis.py} &
    Hausdorff dimension, IFS, multifractal spectra &
    \texttt{box\_counting\_dim()}, \texttt{multifractal\_spectrum()} \\
\texttt{topology\_bridge.py} &
    Topological invariants connecting algebra to geometry &
    \texttt{euler\_characteristic()}, \texttt{betti\_numbers()} \\
\texttt{genesis\_harmonics.py} &
    Harmonic analysis on exceptional groups &
    \texttt{spherical\_harmonics\_e8()}, \texttt{zonal\_polynomial()} \\
\texttt{quantum\_encoding.py} &
    Five quantum encoding schemes for $\EE{7}$ states &
    \texttt{encode\_root()}, \texttt{decode\_state()} \\
\texttt{quantum\_e7\_circuits.py} &
    $\EE{7}$ quantum circuits for IBM Kyiv &
    \texttt{build\_e7\_circuit()}, \texttt{weyl\_gate()} \\
\texttt{quantum\_e8\_circuits.py} &
    $\EE{8}$ quantum circuits (248-dimensional) &
    \texttt{build\_e8\_circuit()} \\
\texttt{quantum\_simulation.py} &
    General quantum simulation utilities &
    \texttt{run\_vqe()}, \texttt{measurement\_statistics()} \\
\texttt{quantum\_lattice\_boltzmann.py} &
    Quantum LBM formulation (D2Q9 + spectral filter) &
    \texttt{QuantumLBM}, \texttt{step()} \\
\texttt{accelerated\_lbm.py} &
    JAX-accelerated LBM on GPU/TPU &
    \texttt{AcceleratedLBM}, \texttt{run\_simulation()} \\
\texttt{jax\_quantum.py} &
    JAX-based quantum circuit evaluation &
    \texttt{jax\_statevector()}, \texttt{jax\_density\_matrix()} \\
\texttt{aqgm\_framework.py} &
    Algebraic quantum-gravity module interface &
    \texttt{AQGMFramework}, \texttt{compute\_amplitudes()} \\
\texttt{unified\_physics.py} &
    Planck units, dimensional analysis, superforce &
    \texttt{planck\_units()}, \texttt{dimensional\_check()} \\
\texttt{inverse\_design.py} &
    Inverse design of metamaterials from lattice data &
    \texttt{design\_from\_lattice()}, \texttt{optimize\_coupling()} \\
\texttt{algebras\_explorer.py} &
    Backend data extraction for interactive explorer &
    \texttt{get\_roots()}, \texttt{get\_adjacency()} \\
\texttt{generate\_interactive\_explorer.py} &
    Generates \texttt{figures/lie\_algebras\_explorer.html} &
    \texttt{build\_explorer()} \\
\texttt{interactive\_explorer\_template.py} &
    WebGL/Three.js HTML template engine &
    \texttt{render\_template()} \\
\texttt{generate\_scientific\_figures.py} &
    Matplotlib + Plotly figure generation &
    \texttt{generate\_all\_figures()} \\
\texttt{generate\_anchor\_diagrams.py} &
    Dynkin diagram generation (TikZ export) &
    \texttt{dynkin\_tikz()} \\
\texttt{data\_handler.py} &
    Manifest loading, SHA-256 checksums, provenance &
    \texttt{load\_manifest()}, \texttt{sha256\_file()} \\
\texttt{config.py} &
    Global configuration and environment variables &
    \texttt{Config}, \texttt{load\_config()} \\
\texttt{optional\_deps.py} &
    Optional dependency detection (JAX, Qiskit) &
    \texttt{has\_jax()}, \texttt{has\_qiskit()} \\
\texttt{viz.py} &
    Common visualisation utilities &
    \texttt{plot\_root\_system()}, \texttt{save\_figure()} \\
\texttt{main.py} &
    CLI entry point and pipeline orchestrator &
    \texttt{main()} \\
\bottomrule
\caption{Module catalog for \texttt{src/mathphysics/}.}
\label{tab:module_catalog}
\end{longtable}

% ==============================================================================
\section{Core Algorithm Catalog}
\label{app:algorithms}
% ==============================================================================

\subsection{Root System Generation}

\begin{algorithm}[H]
\caption{Generate Root System via Positive Root Enumeration}
\label{alg:roots}
\begin{algorithmic}[1]
\Require Cartan matrix $A \in \ZZ^{r \times r}$, simple roots
    $\{\alpha_1, \ldots, \alpha_r\}$
\Ensure Set $\Phi$ of all roots
\State $\Phi^+ \gets \{\alpha_1, \ldots, \alpha_r\}$; queue $Q
    \gets \Phi^+$
\While{$Q \ne \emptyset$}
    \State $\alpha \gets \text{dequeue}(Q)$
    \For{$i = 1, \ldots, r$}
        \State $p \gets \max\{k \ge 0 : \alpha - k\alpha_i \in \Phi^+\}$
        \State $q \gets -\langle\alpha, \alpha_i^\vee\rangle + p$
            \Comment{String length from Cartan}
        \For{$k = 1, \ldots, q$}
            \State $\beta \gets \alpha + k\alpha_i$
            \If{$\beta \notin \Phi^+$}
                \State $\Phi^+ \gets \Phi^+ \cup \{\beta\}$;
                    enqueue$(Q, \beta)$
            \EndIf
        \EndFor
    \EndFor
\EndWhile
\State \Return $\Phi^+ \cup (-\Phi^+)$
\end{algorithmic}
\end{algorithm}

\textbf{Complexity}: $\mathcal{O}(|\Phi| \cdot r)$ time,
$\mathcal{O}(|\Phi|)$ space.  For $\EE{8}$: $|\Phi| = 240$, $r = 8$,
so $< 2000$ operations.

\subsection{Weyl Group Order Computation}

\begin{algorithm}[H]
\caption{Weyl Group Order via Exponent Formula}
\label{alg:weyl}
\begin{algorithmic}[1]
\Require Lie algebra exponents $\{m_1, \ldots, m_r\}$
\Ensure $|W|$
\State $|W| \gets \prod_{i=1}^{r} (m_i + 1)$
    \Comment{Product of $(m_i + 1)$ for each exponent}
\State \Return $|W|$
\end{algorithmic}
\end{algorithm}

Exponents for exceptional algebras:
\begin{equation}
\begin{array}{c|l}
\text{Algebra} & \text{Exponents} \\
\hline
\EE{6} & 1, 4, 5, 7, 8, 11 \\
\EE{7} & 1, 5, 7, 9, 11, 13, 17 \\
\EE{8} & 1, 7, 11, 13, 17, 19, 23, 29 \\
\end{array}
\end{equation}

\subsection{Quartic $\EE{7}$ Invariant}

The quartic invariant $I_4$ governs black hole entropy in
$\mathcal{N}=8$ SUGRA:

\begin{algorithm}[H]
\caption{Quartic $\EE{7}$ Invariant $I_4(p,q)$}
\label{alg:e7_invariant}
\begin{algorithmic}[1]
\Require Charge vector $(p, q) \in \ZZ^{28} \times \ZZ^{28}$
    (56-dimensional)
\Ensure Quartic $\EE{7}$ invariant $I_4$
\State Reshape $(p,q)$ into symplectic vector $v \in \RR^{56}$
\State Compute $\Omega(v,v) = \sum_{i=1}^{28} (p_i q^i - q_i p^i)$
    \Comment{Symplectic form}
\State Build $4 \times 4$ tensor $T_{abcd}$ from $v$ via $\EE{7}$
    structure constants
\State $I_4 \gets \frac{1}{4!}\sum_{a,b,c,d} T_{abcd}\, v^a v^b v^c v^d$
\State \Return $I_4$
\end{algorithmic}
\end{algorithm}

\textbf{Complexity}: $\mathcal{O}(56^4)$ naive; $\mathcal{O}(56^2)$ with
sparsity exploitation ($|$nonzero structure constants$| \ll 56^4$).

\subsection{Fractal Dimension (Box-Counting)}

\begin{algorithm}[H]
\caption{Box-Counting Hausdorff Dimension}
\label{alg:box_counting}
\begin{algorithmic}[1]
\Require Point set $S \subset \RR^d$, box sizes
    $\{\epsilon_1 > \epsilon_2 > \cdots > \epsilon_k\}$
\Ensure Estimated $\hausdim$
\For{$i = 1, \ldots, k$}
    \State Cover $S$ with boxes of side $\epsilon_i$
    \State $N(\epsilon_i) \gets$ number of non-empty boxes
\EndFor
\State Fit $\log N(\epsilon) \sim -\hausdim \log \epsilon$ by linear
    regression on $({\log \epsilon_i}, {\log N(\epsilon_i)})$
\State \Return slope magnitude as $\hausdim$
\end{algorithmic}
\end{algorithm}

\textbf{Complexity}: $\mathcal{O}(k \cdot |S| \cdot \epsilon_{\min}^{-d})$
worst case; GPU-accelerated for large point sets via JAX.

\subsection{Lattice Boltzmann Method (D2Q9 + Spectral Filter)}

\begin{algorithm}[H]
\caption{LBM Time Step with Order-7 Spectral Filter}
\label{alg:lbm}
\begin{algorithmic}[1]
\Require Distribution functions $f_i(x,t)$, relaxation $\tau$,
    filter coefficients $\sigma_k$ ($k=0,\ldots,7$)
\Ensure Updated $f_i(x, t+1)$
\State \textbf{Collision:}
    $f_i^*(x) \gets f_i(x) - \frac{1}{\tau}(f_i(x) - f_i^{\mathrm{eq}}(x))$
\State \textbf{Streaming:}
    $\tilde{f}_i(x + e_i) \gets f_i^*(x)$
\State \textbf{Spectral filter} (in Fourier space):
\State $\quad \hat{f}_i(k) \gets \mathcal{F}[\tilde{f}_i](k)$
\State $\quad \hat{f}_i(k) \gets \hat{f}_i(k) \cdot \sigma(k/k_{\max})$
    \Comment{$\sigma(x) = \exp(-36 x^{36})$ for order-7 filter}
\State $\quad f_i(x, t+1) \gets \mathcal{F}^{-1}[\hat{f}_i](x)$
\State Apply boundary conditions
\State \Return $f_i(x, t+1)$
\end{algorithmic}
\end{algorithm}

\textbf{Complexity per step}: $\mathcal{O}(N^2 \log N)$ for FFT on
$N \times N$ grid.  GPU-accelerated via JAX with
\texttt{jax.numpy.fft.fft2}.

% ==============================================================================
\section{Hardware Requirements and Performance}
\label{app:hardware}
% ==============================================================================

\begin{table}[ht]
\centering
\begin{tabular}{lcccc}
\toprule
\textbf{Task} & \textbf{CPU} & \textbf{RAM} & \textbf{GPU} &
\textbf{Typical Runtime} \\
\midrule
$\EE{7}$ root generation & Any & 1~GB & Not required & $< 1$~s \\
$\EE{8}$ root generation & Any & 1~GB & Not required & $< 1$~s \\
Weyl group enumeration & 4+ cores & 4~GB & Recommended & 2--10~min \\
LBM $256^2$, 1000 steps & 8+ cores & 8~GB & Optional & 5--30~min \\
LBM $512^2$, 10000 steps & 8+ cores & 16~GB & RTX 3080+ & 1--4~h \\
Explorer HTML generation & Any & 2~GB & Not required & $< 30$~s \\
Full figure suite & 4+ cores & 4~GB & Not required & $< 5$~min \\
Full test suite & 4+ cores & 4~GB & Not required & $< 2$~min \\
\bottomrule
\end{tabular}
\caption{Hardware requirements and expected runtimes for key computational tasks.}
\label{tab:hardware_requirements}
\end{table}

\subsection{GPU Acceleration}

JAX-based modules (\texttt{accelerated\_lbm.py}, \texttt{jax\_quantum.py})
detect CUDA availability via \texttt{optional\_deps.has\_jax()}.  When a
GPU is present, JAX JIT-compiles the inner loop; fallback to CPU NumPy is
automatic.

\begin{verbatim}
# Enable GPU (if available):
XLA_PYTHON_CLIENT_PREALLOCATE=false python -m mathphysics.main --gpu
\end{verbatim}

% ==============================================================================
\section{Interactive Algebraic Explorer}
\label{app:explorer}
% ==============================================================================

The interactive explorer \texttt{figures/lie\_algebras\_explorer.html} is
generated by \texttt{make figures}.

\subsection{Architecture}

\begin{itemize}
    \item \textbf{Three.js}: WebGL-based 3D rendering of root systems.
    \item \textbf{Vue.js 3}: Reactive UI for algebra selection and layer
          toggling.
    \item \textbf{Scikit-learn PCA}: Projects high-dimensional roots to 3D
          preserving maximal variance.
    \item \textbf{Tailwind CSS}: Responsive control panel.
\end{itemize}

\subsection{Supported Algebras}

Real-time switching between: $\EE{4}$ through $\EE{11}$, $F_4$, $G_2$,
$D_5$, and $A_n$ for $n \le 8$.

\subsection{Layer Controls}

\begin{enumerate}
    \item \textbf{Root Layer}: 3D position of each root vector.
    \item \textbf{Weight/Connection Layer}: Nearest-neighbour edges from
          Cartan matrix.
    \item \textbf{Clifford Mapping}: Colour-coded Cayley--Dickson bit
          patterns.
    \item \textbf{Metadata Tooltips}: Coordinates and representation labels
          on hover.
\end{enumerate}

\subsection{Accessibility}

The explorer adheres to WCAG~2.1 AA: keyboard navigation for all
controls, screen-reader ARIA labels on all interactive elements, high-contrast
mode toggle.

% ==============================================================================
\section{Simulation Engine (LBM)}
\label{app:lbm_engine}
% ==============================================================================

Physical simulations use a JAX-accelerated Lattice Boltzmann engine:

\begin{itemize}
    \item \textbf{Velocity set}: D2Q9 (two-dimensional, nine velocities).
    \item \textbf{Equilibrium}: Maxwell--Boltzmann to second order in $u/c_s$.
    \item \textbf{Filter}: Order-7 exponential spectral filter
          $\sigma(\xi) = \exp(-36\xi^{36})$ following~\cite{hou2007computing}.
    \item \textbf{GPU offload}: CUDA via JAX; automatic CPU fallback.
    \item \textbf{Output}: Parquet snapshots every $N_{\mathrm{snap}}$ steps;
          metadata including step, timestamp, and SHA-256 stored in
          \texttt{data/registry/simulation\_registry.toml}.
\end{itemize}

\subsection{Validation}

LBM correctness is verified by:
\begin{enumerate}
    \item Mass conservation: $\sum_x \rho(x,t) = \text{const}$ to
          machine precision.
    \item Momentum conservation: $\sum_x \rho(x,t) u(x,t) =
          \text{const}$ (periodic BC).
    \item Poiseuille flow: analytical solution recovered to $< 0.1\%$
          error at Re $= 100$.
    \item Rhines scale: $L_R = \sqrt{U_\star / \beta}$ reproduced to
          $< 5\%$ in jet-forming runs.
\end{enumerate}

% ==============================================================================
\section{Build Pipeline}
\label{app:build}
% ==============================================================================

\subsection{Makefile Targets}

\begin{verbatim}
make lint          # ruff + flake8; treat warnings as errors
make check-types   # mypy --strict
make test          # pytest with PYTHONHASHSEED=0
make benchmark     # timing suite for all algorithms
make figures       # generate all figures + explorer HTML
make papers        # pdflatex + bibtex + pdflatex x2
make fetch-all     # download all bibliography PDFs
make verify-checksums # SHA-256 integrity check
make cleanbuild    # clean + all of the above in sequence
\end{verbatim}

\subsection{CI/CD}

The GitHub Actions workflow \texttt{.github/workflows/ci.yml} runs:
lint, type check, tests, and paper compilation on every push to
\texttt{main}.  All quality gates are hard failures (no
\texttt{continue-on-error}).
